
\vspace{0.5cm}
\begin{tikzpicture}[font=\footnotesize]

% Start block
\node[draw,
    rounded corners,
    rectangle,
    align=center,
    text width=2.1cm,
    minimum width=2.2cm,
    minimum height=1cm] (block0) 
    {
    Adjust by \\
    $t_\mathrm{cal.}+\SI{5}{\second}$ \\
    $\delta-0.01\,\mathrm{ADC}$
    };

\node[draw,
    align=center,
    text width=4.2cm,
    above=of block0,
    minimum width=3cm,
    minimum height=1cm
    ] (block1) {
    $\Delta r > 0:\,T_{70}^\mathrm{new} = T_{70}^\mathrm{old} + \delta$
    $\Delta r < 0:\,T_{70}^\mathrm{new} = T_{70}^\mathrm{old} - \delta$
    };

\node[draw,
    shape=diamond,
    left=0.6cm of block1,
    minimum height=1cm
    ] (block3) { 
    Wait $t_\mathrm{cal.}$
    };

\node[draw,
    align=center,
    minimum height=1cm,
    below=of block0
    ] (block4) {
    $\Delta r < \sigma_r$
    };

\node[draw,
    align=center,
    left=1.54cm of block4,
    minimum width=3cm,
    minimum height=1cm
] (block5) {
	$\Delta r = \left|f_{70} - \SI{70}{\hertz}\right|$ \\ $\sigma_{r} = \sqrt{n\,/\,t_\mathrm{cal.}}$
    };

\node[draw,
    rounded corners,
    rectangle,
    right=0.4cm of block0,
    text width=2.1cm,
    align=center,
    minimum width=2.2cm,
    minimum height=1cm] (block6) 
    {
    Set increment\\
$\delta = 0.1\,\mathrm{ADC}.$
    };

\node[draw,
    align=center,
    minimum height=1cm,
    below=1.017cm of block6
    ] (block7) {
    $\Delta r < \SI{5}{\hertz}$
    };

\node[draw,
    rounded corners,
    rectangle,
    right=0.4cm of block6,
    text width=2.1cm,
    align=center,
    minimum height=1cm] (block8) 
    {
    Leave $t_\mathrm{cal.}$ and $\delta$ invariant
    };

\node[draw,
    align=center,
    minimum height=1cm,
    below=0.968cm of block8
    ] (block9) {
    $\Delta r < \SI{20}{\hertz}$
    };

\node[draw,
    rounded corners,
    rectangle,
    right=0.4cm of block8,
    text width=2.1cm,
    align=center,
    minimum height=1cm] (block10) 
    {
    Set values \\
    $t_\mathrm{cal.} = \SI{10}{\second}$, \\
    $\delta = 1.0\,\mathrm{ADC}$.
    };

\node[draw,
    rounded rectangle,
    above=3.55cm of block8,
    align=center,
    minimum width=2.2cm,
    minimum height=1cm] (block13) 
    {START};

\node[draw,
    rounded rectangle,
    left=1.9cm of block13,
    align=center,
    minimum width=2.2cm,
    minimum height=1cm] (block2) 
    {
    Initialize $T_{70} = 50\,\mathrm{ADC}$, $t_\mathrm{cal.} = \SI{10}{\second}$, and $\delta = 0.2\,\mathrm{ADC}.$
    };

\node[coordinate,right=0.54cm of block1] (block11) {};
\node[coordinate,right=2.764cm of block11] (block12) {};

% Arrows
\draw[-latex] (block2) edge (block3);
\draw[-latex] (block3) edge (block5);
\draw[-latex] (block5) edge (block4);
\draw[-latex] (block13) edge (block2);

\draw[-latex] (block4) edge node[pos=0.45,fill=white,inner sep=2pt]{Yes}(block0);
\draw[-latex] (block4) edge node[pos=0.45,fill=white,inner sep=2pt]{x}(block7);
\draw[-latex] (block7) edge node[pos=0.4,fill=white,inner sep=2pt]{Yes}(block6);
\draw[-latex] (block7) edge node[pos=0.45,fill=white,inner sep=2pt]{x}(block9);
\draw[-latex] (block9) edge node[pos=0.4,fill=white,inner sep=2pt]{Yes}(block8);
    
\draw[-latex] (block9) -| (block10)
    node[pos=0.15,fill=white,inner sep=1]{No};

\draw[-latex] (block1) edge (block3);
\draw[-latex] (block0) edge (block1);
\draw (block8) edge (block12);
\draw (block6) -- (block11);
\draw[-latex] (block10) |- (block1);

\end{tikzpicture}
\vspace{0.5cm}
